\uuid{8k2D}
\titre{ Diagonalisation et calcul de puissance }

\niveau{L2}
\module{ Réduction et diagonalisation }
\chapitre{Application linéaire}
\sousChapitre{Diagonalisation}
\theme{}
\auteur{Q. Liard}
\datecreate{ 2026-03-05}
\organisation{AMSCC}
\difficulte{}




\contenu{

\texte{
On considère la matrice
\[
A=
\begin{pmatrix}
4 & 0 \\
1 & 3
\end{pmatrix}.
\]
}

\begin{enumerate}

\item \question{Calculer le polynôme caractéristique de $A$.}
\reponse{$\chi_A(X)=(X-4)(X-3)$.}

\item \question{Déterminer les valeurs propres et les vecteurs propres de $A$.}
\reponse{Valeurs propres : $4$ et $3$.  
Vecteurs propres associés respectivement : $(1,1)$ et $(0,1)$.}

\item \question{Montrer que $A$ est diagonalisable.}
\reponse{Deux valeurs propres distinctes donc diagonalisable.}

\item \question{En déduire l’expression de $A^n$.}
\reponse{Si $A=PDP^{-1}$ alors $A^n=PD^nP^{-1}$ avec $D^n=\mathrm{diag}(4^n,3^n)$.}

\end{enumerate}

}